\chapter{Before the First Feature}

The screen backend from the last chapter holds a session with the plant, keeps
the last value of every node and pushes changes to the browsers. Before it can
do any of that, it has to exist as a program: a process that starts, listens on
a port and answers.

It's written in JavaScript, running on Node, because there is a mature OPC UA
client for Node and because the browser on the other side speaks the same
language. A Node server in 2026 is assembled from a dozen tools that weren't
there in 2009: a compiler, a module system, a lock file, a framework, a schema
library, a linter, a formatter, a test runner. Each one has a history, and
almost all of them exist for the same reason: to make a mistake show up
earlier.

\begin{figure}[htbp]
\centering
\begin{tikzpicture}[x=1.8cm, y=-0.78cm]
  \foreach \x/\s in {0/lint, 1/typecheck, 2/test, 3/startup, 4/a request, 5/someone\\notices} {
    \node[hf/note, anchor=south, font=\sffamily\footnotesize\bfseries] at (\x,0.2) {\s};
    \draw[hf/lifeline] (\x,0.35) -- (\x,4.6);
  }
  \foreach \y/\t/\from/\to in {
      1/a misspelled method/4/1,
      2/a promise nobody waits for/5/0,
      3/an invalid setting/4/3,
      4/a route that stops answering/5/2} {
    \node[hf/note, anchor=east, align=right] at (-0.25,\y) {\t};
    \fill[hfAlert] (\from,\y) circle (2.6pt);
    \fill[hfReader] (\to,\y) circle (2.6pt);
    \draw[hf/ask, shorten >=3pt, shorten <=3pt] (\from,\y) -- (\to,\y);
  }
\end{tikzpicture}
\caption{Where a mistake shows up. Red: with only Node and a text editor. Blue:
after the tools in this chapter. Each tool moves one kind of mistake to the
left, closer to the person who made it.}
\end{figure}

\section{What Node is}

\begin{history}{1999--2009 \textperiodcentered{} Ten thousand connections}
In 1999, Dan Kegel published ``The C10K problem'': how to make one server handle
ten thousand clients at once. The standard design of the time, Apache's, gave
each connection its own process or thread. A thread costs memory even when it's
doing nothing, and most of a web connection's life is doing nothing: waiting
for the network, the disk, a database. Ten thousand connections meant ten
thousand threads, almost all of them asleep.

The answer was to let one thread wait on all the connections at once and only
act when one of them had something. Operating systems added the calls for it
(\texttt{epoll} on Linux in 2002, \texttt{kqueue} on BSD in 2000), and in 2004
Igor Sysoev released nginx, a web server built around one loop.

Writing applications that way was still painful, because every library in
every language assumed it could block. In 2008 Google released Chrome with V8,
a fast JavaScript engine from Lars Bak's team in Denmark. Ryan Dahl noticed
that JavaScript was the one popular language with no blocking I/O libraries at
all: in the browser, everything was already a callback. In November 2009, at
JSConf EU in Berlin, he presented Node: V8 plus an event loop plus a set of
libraries in which nothing blocks. He had started from a small problem: an
upload progress bar, which is hard to build on a server that can't report
anything while it's still receiving the file.
\end{history}

\begin{concept}{Event loop}
One thread runs all the program's code. When the code needs to wait (for a
socket, a file, a timer), it registers a callback and returns. The loop asks
the operating system which waits are over and runs the callbacks for those,
one at a time.

\textbf{What it fixes:} a thousand idle connections cost a thousand small
objects, not a thousand threads.

\textbf{What it costs:} while one callback is computing, nobody else runs. A
loop that takes two seconds to parse a file freezes every connection for two
seconds.
\end{concept}

\begin{figure}[htbp]
\centering
\begin{tikzpicture}[x=1mm, y=1mm]
  % threads
  \node[hf/note, anchor=south west, font=\sffamily\footnotesize\bfseries] at (0,30) {one thread per connection};
  \foreach \r in {0,1,2,3} {
    \pgfmathsetmacro{\y}{24-7*\r}
    \fill[black!8] (0,\y) rectangle (62,\y+4);
    \node[hf/note, anchor=east] at (-1,\y+2) {thread \the\numexpr\r+1\relax};
  }
  \foreach \x/\r in {2/0, 30/0, 9/1, 44/1, 16/2, 51/2, 23/3, 37/3} {
    \pgfmathsetmacro{\y}{24-7*\r}
    \fill[hfReader] (\x,\y) rectangle (\x+4,\y+4);
  }
  \node[hf/note] at (31,-1) {mostly asleep, each with its own memory};
  % loop
  \begin{scope}[xshift=72mm]
    \node[hf/note, anchor=south west, font=\sffamily\footnotesize\bfseries] at (0,30) {one event loop};
    \fill[black!8] (0,13.5) rectangle (62,17.5);
    \node[hf/note, anchor=east] at (-1,15.5) {loop};
    \foreach \x/\n in {2/1, 9/2, 16/3, 23/4, 30/1, 37/4, 44/2, 51/3} {
      \fill[hfReader] (\x,13.5) rectangle (\x+4,17.5);
      \node[font=\sffamily\tiny, text=white] at (\x+2,15.5) {\n};
    }
    \node[hf/note] at (31,6) {the same work, one thread;\\the OS watches the sockets};
  \end{scope}
\end{tikzpicture}
\caption{Four connections, the same work. On the left, four threads that spend
most of their time waiting. On the right, one loop that runs each piece of work
when its data arrives.}
\end{figure}

Dahl's description, V8 plus an event loop plus libraries in which nothing
blocks, names Node's three parts. None of them was new, and each came from
somewhere else.

The first is the \textbf{engine}. The JavaScript language, as its specification
defines it, has no input or output at all: no files, no network, not even
\texttt{setTimeout}. An engine like V8 runs the language and nothing more. It
turns your code into machine instructions, keeps the objects your code creates
in a region of memory called the \emph{heap}, and keeps track of which function
is running on the \emph{call stack}: each call puts an entry on top, each return
takes one off. Whatever program hosts the engine has to supply the rest. For
most of its first fifteen years, the host that mattered was the browser, which
supplied the page, the clicks and the network requests. V8 was written to be
hosted by other programs too: Google published it in 2008 as a C++ library any
program could embed, and that's what Dahl did.

The second is the \textbf{loop}, and Dahl didn't write his own. Node~0.1 used
libev, an event loop for Unix that Marc Lehmann had published in 2007, which
waits on sockets with \texttt{epoll} or \texttt{kqueue}, the calls from the
C10K story. Files were a different problem. Linux had no way to ask for a read
of an ordinary file and be told later when it was done, so Node used another of
Lehmann's libraries, libeio, which does each read on a small pool of threads:
the thread blocks on the read so the loop doesn't have to. In 2011, Joyent and
Microsoft ported Node to Windows, where neither library ran. The port produced
libuv, which at first wrapped libev and libeio on Unix and, on Windows, IOCP,
the system's own way of waiting on many sockets. Later it replaced Lehmann's two
libraries with its own code. libuv kept the same
split: sockets go to the operating system, and whatever can only block goes to
a pool of four threads.

The pool is the part most often drawn wrong. The diagram of Node that circulates
most on the web puts \emph{network} among the worker threads, and that's a
mistake. Sockets never touch the pool: the loop hands them to the operating
system, which watches all of them for one thread. The pool is for the work that
has no such call on every system: reading files, some cryptography and
compression, and looking a name up in DNS, whose system call, \texttt{getaddrinfo},
blocks. The difference is easy to see. With the four pool threads busy
computing four slow password hashes (\texttt{crypto.pbkdf2}, which runs on the
pool), a request to a server on the same machine took 22 milliseconds when the
address was \texttt{127.0.0.1} and 1.7 seconds when it was \texttt{localhost}.
The number needs no lookup and goes straight to the operating system. The name
had to wait in line for a free thread, and one hash takes 1.5 seconds.

The third part is what lets JavaScript reach the other two. V8 and libuv are
libraries written in C and C++, and JavaScript can't call a C function by
itself. The connection is a \textbf{binding}: a C++ function registered in the
engine under a JavaScript name, so that when your code calls that name, the
engine runs the C++ function instead. The idea is as old as scripting
languages. John Ousterhout designed Tcl in 1988 to be embedded in C programs,
and each program added commands of its own, written in C, that Tcl scripts
could call by name. Node has a binding for each thing its libraries need from
below: open a socket, read a file, start a timer. On top of the bindings sits
Node's standard library (\texttt{node:http}, \texttt{node:fs},
\texttt{node:net}), ordinary JavaScript that gives those calls a friendlier
shape. Figure~\ref{fig:node-layers} puts the parts together.

\begin{figure}[htbp]
\centering
\begin{tikzpicture}[x=1cm, y=1cm]
  \tikzset{
    layer/.style={draw=hfBus, fill=hfBus!6, rounded corners=3pt, line width=0.8pt},
    part/.style={hf/base, draw=hfBus, fill=white, minimum width=24mm, font=\sffamily\footnotesize},
    small/.style={hf/base, draw=hfReader, dashed, fill=hfReader!5, minimum width=24mm, font=\sffamily\footnotesize},
    ph/.style={hf/base, draw=hfBus, fill=white, font=\sffamily\scriptsize, inner sep=2pt, minimum width=12mm},
    both/.style={{Stealth[length=2.2mm]}-{Stealth[length=2.2mm]}, line width=0.8pt, black!55},
  }
  % your code and V8
  \node[hf/reader, minimum width=30mm, minimum height=9mm] (app) at (0,3.45) {your code\\{\footnotesize JavaScript}};
  \draw[layer] (-1.5,-1.6) rectangle (1.5,2.55);
  \node[font=\sffamily\bfseries] at (0,2.2) {V8};
  \node[hf/note] at (0,1.8) {runs JavaScript};
  \node[part] at (0,1.1) {call stack};
  \node[part] at (0,0.35) {heap};
  \draw[both] (app.south) -- (0,2.55);
  % Node
  \draw[layer] (2.0,-1.6) rectangle (4.8,3.9);
  \node[font=\sffamily\bfseries] at (3.4,3.55) {Node};
  \node[part] at (3.4,2.55) {standard library\\{\scriptsize \texttt{node:http}, \texttt{node:fs}}\\{\scriptsize JavaScript}};
  \node[part] (bind) at (3.4,-0.45) {bindings\\{\scriptsize C++}};
  \draw[both] (1.5,-0.45) -- (bind.west);
  \draw[both] (3.4,1.95) -- (bind.north);
  % libuv
  \draw[layer] (5.3,-1.6) rectangle (12.4,3.9);
  \node[font=\sffamily\bfseries] at (8.85,3.55) {libuv};
  \draw[both] (bind.east) -- (5.3,-0.45);
  \node[font=\sffamily\footnotesize\bfseries, text=hfBus] at (7.6,1.0) {event loop};
  \draw[hfBus, line width=1.4pt, -{Stealth[length=3mm]}] (7.6,2.45) arc[start angle=90, end angle=-225, radius=1.45];
  \coordinate (pol) at (8.6,-0.05);
  \node[hf/note, font=\sffamily\footnotesize\bfseries] at (11.15,3.0) {thread pool};
  \foreach \y/\n in {2.4/1, 1.75/2, 1.1/3, 0.45/4}
    \node[hf/source, minimum width=16mm, minimum height=4.5mm, font=\sffamily\scriptsize, inner sep=1pt] (t\n) at (11.15,\y) {thread \n};
  % operating system
  \node[hf/source, minimum width=103mm, minimum height=8mm] (os) at (7.2,-2.55)
    {operating system \quad {\footnotesize \texttt{epoll} (Linux), \texttt{kqueue} (macOS, BSD), IOCP (Windows)}};
  \draw[both] (pol) -- (pol |- os.north) node[midway, right, hf/note] {sockets};
  \draw[both] (t4.south) -- (t4.south |- os.north) node[midway, left, hf/note, align=right] {files, DNS,\\crypto};
\end{tikzpicture}
\caption{Where Node sits. V8 runs the JavaScript, yours and Node's standard
library; the bindings connect it to libuv, which owns the event loop and a pool
of four threads. The operating system does the waiting: it watches the sockets
for the loop and runs the blocking calls for the pool.}
\label{fig:node-layers}
\end{figure}

That combination, an engine with what it needs to touch the world, is a
\emph{runtime}.

Node is a runtime, and the difference between a runtime and a framework comes
down to who calls whom.

\begin{history}{1985--1988 \textperiodcentered{} Don't call us, we'll call you}
In 1985, Richard Sweet described the Mesa programming environment at Xerox
PARC and named one of its rules ``Hollywood's Law'': \emph{don't call us,
we'll call you}. The user interface code called the application when something
happened, never the other way around. In 1988, Ralph Johnson and Brian Foote
used that inversion to define a \textbf{framework} in ``Designing Reusable
Classes'': a set of classes that make up a design for a family of programs,
where the framework holds the flow of control and calls the code you plug into
it.
\end{history}

\begin{concept}{Runtime, library, framework}
A \emph{runtime} runs your program and gives it access to the machine: memory,
time, files, network. A \emph{library} is code you call: you decide when, and
it returns. A \emph{framework} is code that calls you: it owns the flow, and you
fill in the parts it leaves open, like what to do when a request arrives on a
given route.

\textbf{What it fixes:} the framework already made the decisions every program
of that kind needs (routing, parsing, errors), so you only write the part that
is different.

\textbf{What it costs:} you write your program the framework's way, and leaving
it later means rewriting more than a few calls.
\end{concept}

Node gives you the runtime and some libraries, but no framework. Its
\texttt{http} module parses the bytes of a request and hands you the result,
and everything after that is yours to write.

\begin{code}
import http from 'node:http';

http.createServer((req, res) => {
  if (req.method === 'GET' && req.url === '/status') {
    res.setHeader('content-type', 'application/json');
    res.end(JSON.stringify({ ok: true }));
    return;
  }
  res.statusCode = 404;
  res.end();
}).listen(8080);
\end{code}

That works until a client adds anything to the address:

\begin{terminal}
GET  /status        -> 200
GET  /status?x=1    -> 404
POST /status        -> 404
\end{terminal}

\texttt{req.url} is the raw path with the query string attached, so comparing
it as a string fails as soon as someone adds \texttt{?x=1}. Fixing that means
parsing the URL. Answering the wrong method properly means another branch.
Reading a body means collecting chunks from a stream and parsing them. Writing
a server on the bare \texttt{http} module is a good way to learn what a request
is, and a mistake past the first route.

\section{What runs next}

The event loop in Figure~\ref{fig:node-layers} runs one callback at a time. What
decides which one comes next? The answer grew over ten years, in the browser and
in Node at the same time, and its names only make sense in the order they
appeared.

\begin{history}{2008 \textperiodcentered{} A queue of tasks}
Browsers had always run one thing at a time and kept the rest waiting, but no
specification said in what order. In August 2008, Ian Hickson, editing the HTML5
draft at the WHATWG, wrote it down: each page has an \emph{event loop}, and
everything that happens later (a click, a timer, a finished download) goes into
a \emph{task queue} as a \emph{task}. The loop takes one task, runs it to the
end, and takes the next.
\end{history}

Node's loop works the same way, with more structure. libuv goes around in
\emph{turns}, and each turn passes through the same stops, in the same order,
called \emph{phases}. Each phase has its own queue, fed by one kind of event:

\begin{itemize}
\item \textbf{Timers} runs the callbacks of \texttt{setTimeout} and
  \texttt{setInterval} whose time has come.
\item \textbf{Pending} runs a few I/O callbacks the operating system left for
  the next turn, such as some network errors. It rarely matters.
\item \textbf{Poll} is the main stop. The loop asks the operating system which
  sockets have data and which file reads have finished, and runs their
  callbacks: a request coming in, a message from the plant. When no other
  queue has anything, the loop sleeps here until something arrives or the next
  timer is due.
\item \textbf{Check} runs, right after poll, the callbacks scheduled for that
  moment. Their story comes below.
\item \textbf{Close} runs the callbacks of sockets and other handles that have
  just closed.
\end{itemize}

Then the loop goes back to timers, and that is one turn. Each callback the loop
runs from a phase is a task in Hickson's sense, though Node's documentation
rarely uses the word.

Since January 2010, Node has also had a way to run something \emph{between}
two tasks. An object that announces events, like a server or a stream, has a
problem when something goes wrong while it's still being created. If it
announces the error right away, nobody is listening yet: the code that created
it hasn't finished, so it hasn't attached a listener. \texttt{process.nextTick}
solves that. The function you give it runs as soon as the current code
finishes, before the loop moves on to anything else, and by then the listener
is there. It's a second, small queue, and Node empties it after every task.

The opposite wish, to run something only after the loop has looked at
everything else that's waiting, got its tool from Microsoft. \texttt{setImmediate}
was proposed for browsers in 2011 as ``Efficient Script Yielding'', shipped in
Internet Explorer~10, was turned down by the other browsers, and landed in
Node~0.10 in 2013, where it stayed as the check phase. Node's own documentation
admits that the two names should be swapped: \texttt{nextTick} runs right away,
and \texttt{setImmediate} on the next pass.

The browser reached the same idea as \texttt{nextTick} a few years later, for a
different reason.

\begin{history}{2011--2012 \textperiodcentered{} A queue between tasks}
In 2011, engineers from Google and Mozilla, among them Rafael Weinstein, Adam
Klein and Jonas Sicking, were replacing the DOM's mutation events, which fired
in the middle of every change to the page and made pages slow. Their
replacement, \texttt{MutationObserver}, delivers the changes in a batch, at the
end of the script that made them. One whole task later was too late, because
the browser may draw the page between two tasks. In 2012 the HTML specification
grew a second, small queue for that moment, emptied at a ``microtask
checkpoint'' after each task. Its entries are \emph{microtasks}.
\end{history}

Then came promises. A promise is an object that stands for a result that isn't
ready yet: a function that has to wait returns one at once, and whoever holds it
attaches, with \texttt{.then()}, the code that should run when the result
arrives. Chapter~4 tells where promises came from. Here the only question is
which queue that code goes into, and by 2012 there were two to choose from.

\begin{history}{2012--2018 \textperiodcentered{} Promises take the small queue}
Promise libraries were many and didn't work with each other until their authors
agreed on a short specification, Promises/A+. Version 1.0, in December 2012,
required that a \texttt{then} callback never run before the code that
registered it had finished. In August 2013, Isaac
Schlueter, then Node's lead, explained why in ``Designing APIs for
Asynchrony'': an API that calls its callback sometimes right away and sometimes
later ``releases Zalgo'', and the code that uses it can't be reasoned about.
Either queue met the rule, but a task waits behind every timer and socket in
the loop, so the libraries that cared about speed used the small one.

ES2015 put \texttt{Promise} in the language, and with it the first queue the
language itself defines. The specification calls its entries \emph{jobs}: when
a promise settles, its callbacks become jobs, and jobs run when no other code
is running. The language says nothing more, because it has no event loop; the
host picks the moment. Browsers and Node both run promise jobs as microtasks.
\texttt{async} and \texttt{await} arrived on top of them in Node~7.6, in
February 2017, a few months before ES2017 made them official. With them, a
function can write \texttt{await} in front of a promise and stop there; the rest
of the function becomes the \texttt{.then()} callback, so every line after an
\texttt{await} runs as a microtask.

The last change was in Node~11, in 2018. Until then, Node emptied the microtask
queue only between phases, so a promise created in one timer callback waited
until every other timer due in that turn had run. Browsers empty it after every
callback, and the same code printed in a different order in the two. Node~11
took the browser's rule. With two timers, where the first one resolves a
promise, Node today prints the first timer, then the promise, then the second
timer.
\end{history}

\begin{etymology}
\emph{Micro-} and \emph{macro-} are Greek for small and large. Once there was a
microtask, the old kind needed a name to tell it apart, and ``macrotask'' is
that name, made after the fact, like \emph{acoustic guitar} after the electric
one. In September 2013, version 1.1 of Promises/A+ put the two side by side: a
library could defer a callback with a ``macro-task'' mechanism such as
\texttt{setTimeout} or \texttt{setImmediate}, or with a ``micro-task'' mechanism
such as \texttt{MutationObserver} or \texttt{process.nextTick}. The HTML
specification never took the word. There it's still a task.
\end{etymology}

\begin{figure}[htbp]
\centering
\begin{tikzpicture}[x=1.15cm, y=1cm]
  \draw[black!40, line width=0.8pt] (-0.3,0) -- (10.3,0);
  \foreach \yr/\up/\txt in {
      2008/1/{HTML5 defines the event loop and its tasks},
      2009/0/{Node~0.1 runs on libev},
      2010/1/{\texttt{nextTick} in Node~0.1.26},
      2011/0/{libuv, for Windows; mutation observers},
      2012/1/{microtasks enter HTML; Promises/A+~1.0},
      2013/0/{\texttt{setImmediate} in Node~0.10; Promises/A+~1.1 names the ``macro-\allowbreak{}task''},
      2015/1/{\texttt{Promise} and its jobs in ES2015},
      2017/0/{\texttt{async} and \texttt{await} in Node~7.6},
      2018/1/{Node~11 empties microtasks after every callback}} {
    \pgfmathsetmacro{\x}{\yr-2008}
    \fill[hfBus] (\x,0) circle (2.4pt);
    \node[hf/note, font=\sffamily\scriptsize\bfseries] at (\x,{\up ? -0.3 : 0.3}) {\yr};
    \ifnum\up=1
      \draw[black!25] (\x,0.1) -- (\x,0.55);
      \node[hf/note, anchor=south, text width=22mm, font=\sffamily\scriptsize] at (\x,0.55) {\txt};
    \else
      \draw[black!25] (\x,-0.1) -- (\x,-0.55);
      \node[hf/note, anchor=north, text width=22mm, font=\sffamily\scriptsize] at (\x,-0.55) {\txt};
    \fi
  }
\end{tikzpicture}
\caption{How the loop's queues got their names, from 2008 to 2018.}
\end{figure}

Put together, the names sort themselves. A task belongs to the host: libuv's
phases in Node, the HTML event loop in a browser. A microtask belongs to the
language: the queue that promises use is defined by the language and kept by
the engine, V8, and the host only decides when to empty it. \texttt{process.nextTick} belongs to Node alone.
After every task, Node empties \texttt{nextTick} first, then the microtasks,
including any added along the way, and only then lets the loop move on.

\begin{concept}{Tasks and microtasks}
A \emph{task} (in blog posts, a macrotask) is one callback the host's loop runs
because something outside the program happened: a timer, a socket, a file. A
\emph{microtask} is a callback the program scheduled to run right after the
current one, almost always a promise continuation. After every task, Node runs every
microtask in the queue, including the ones added along the way, and only then
moves on to the next task.

\textbf{What it fixes:} the rest of an \texttt{async} function continues as
soon as its value is ready, without waiting behind timers and sockets, and
always after the current code, never in the middle of it.

\textbf{What it costs:} microtasks can starve the loop. While their queue isn't
empty, no timer fires, no socket is read and no request is answered.
\end{concept}

Four lines are enough to see the order, one for each queue:

\begin{code}
setTimeout(() => console.log('timer'), 0);
Promise.resolve().then(() => console.log('promise'));
process.nextTick(() => console.log('tick'));
console.log('script');
\end{code}

\begin{terminal}
$ node order.js
script
tick
promise
timer
\end{terminal}

The timer was set first, for zero milliseconds, and it came last. The
\texttt{nextTick} was set after the promise, and it came before it.
Figure~\ref{fig:order-frames} shows why, one moment at a time.

\begin{figure}[htbp]
\centering
\begin{tikzpicture}[x=1cm, y=1cm]
  \tikzset{
    fr/.style={hf/base, draw=hfReader, fill=hfReader!10, font=\ttfamily\scriptsize,
      minimum width=23mm, minimum height=4.4mm, inner sep=1pt},
    nq/.style={fr, dashed, draw=hfBus, fill=hfBus!4, minimum width=19mm},
    mq/.style={fr, dashed, fill=hfReader!4, minimum width=19mm},
    tq/.style={fr, draw=hfBus, fill=hfBus!10, minimum width=19mm},
    printed/.style={font=\ttfamily\scriptsize, align=left, anchor=north west},
    when/.style={hf/note, anchor=west, align=left, text width=31mm},
    head/.style={hf/note, font=\sffamily\footnotesize\bfseries, align=center},
    none/.style={hf/note, text=black!35},
  }
  \def\rh{1.5}
  \node[head] at (4.75,0.8) {call stack};
  \node[head, text=hfBus] at (7.05,0.8) {\texttt{nextTick}\\queue};
  \node[head, text=hfReader] at (9.15,0.8) {microtask\\queue};
  \node[head, text=hfBus] at (11.25,0.8) {timers\\(libuv)};
  \node[head] at (13.3,0.8) {printed};
  \foreach \r in {1,...,6} \draw[black!12] (-0.1,{-\rh*\r+0.75}) -- (14.4,{-\rh*\r+0.75});
  % 1
  \node[when] at (0,0) {\textbf{1.} line 1: the timer is handed to libuv};
  \node[fr] at (4.75,-0.23) {order.js};
  \node[fr] at (4.75,0.24) {setTimeout};
  \node[none] at (7.05,0) {empty};
  \node[none] at (9.15,0) {empty};
  \node[tq] at (11.25,0) {'timer'};
  % 2
  \node[when] at (0,{-1*\rh}) {\textbf{2.} line 2: the promise is already resolved};
  \node[fr] at (4.75,{-1*\rh-0.23}) {order.js};
  \node[fr] at (4.75,{-1*\rh+0.24}) {then};
  \node[none] at (7.05,{-1*\rh}) {empty};
  \node[mq] at (9.15,{-1*\rh}) {'promise'};
  \node[tq] at (11.25,{-1*\rh}) {'timer'};
  % 3
  \node[when] at (0,{-2*\rh}) {\textbf{3.} line 3: a callback for Node's own queue};
  \node[fr] at (4.75,{-2*\rh-0.23}) {order.js};
  \node[fr] at (4.75,{-2*\rh+0.24}) {nextTick};
  \node[nq] at (7.05,{-2*\rh}) {'tick'};
  \node[mq] at (9.15,{-2*\rh}) {'promise'};
  \node[tq] at (11.25,{-2*\rh}) {'timer'};
  % 4
  \node[when] at (0,{-3*\rh}) {\textbf{4.} line 4 prints};
  \node[fr] at (4.75,{-3*\rh-0.23}) {order.js};
  \node[fr] at (4.75,{-3*\rh+0.24}) {console.log};
  \node[nq] at (7.05,{-3*\rh}) {'tick'};
  \node[mq] at (9.15,{-3*\rh}) {'promise'};
  \node[tq] at (11.25,{-3*\rh}) {'timer'};
  \node[printed] at (12.55,{-3*\rh+0.2}) {script};
  % 5
  \node[when] at (0,{-4*\rh}) {\textbf{5.} the script ends; Node empties \texttt{nextTick} first};
  \node[fr] at (4.75,{-4*\rh-0.23}) {'tick'};
  \node[fr] at (4.75,{-4*\rh+0.24}) {console.log};
  \node[none] at (7.05,{-4*\rh}) {empty};
  \node[mq] at (9.15,{-4*\rh}) {'promise'};
  \node[tq] at (11.25,{-4*\rh}) {'timer'};
  \node[printed] at (12.55,{-4*\rh+0.2}) {script\\tick};
  % 6
  \node[when] at (0,{-5*\rh}) {\textbf{6.} then the microtasks};
  \node[fr] at (4.75,{-5*\rh-0.23}) {'promise'};
  \node[fr] at (4.75,{-5*\rh+0.24}) {console.log};
  \node[none] at (7.05,{-5*\rh}) {empty};
  \node[none] at (9.15,{-5*\rh}) {empty};
  \node[tq] at (11.25,{-5*\rh}) {'timer'};
  \node[printed] at (12.55,{-5*\rh+0.3}) {script\\tick\\promise};
  % 7
  \node[when] at (0,{-6*\rh}) {\textbf{7.} the loop reaches the timers phase};
  \node[fr] at (4.75,{-6*\rh-0.23}) {'timer'};
  \node[fr] at (4.75,{-6*\rh+0.24}) {console.log};
  \node[none] at (7.05,{-6*\rh}) {empty};
  \node[none] at (9.15,{-6*\rh}) {empty};
  \node[none] at (11.25,{-6*\rh}) {empty};
  \node[printed] at (12.55,{-6*\rh+0.4}) {script\\tick\\promise\\timer};
\end{tikzpicture}
\caption{The four lines, one moment at a time. The call stack grows upward: the
box on top is the function running now. A word in quotes is the callback that
prints it. When the stack empties, Node empties its own \texttt{nextTick} queue
first and the microtask queue second; the timer callback waits in libuv until
the loop comes around to the timers phase.}
\label{fig:order-frames}
\end{figure}

Lines 1 to 4 are one task: running the script. Each line puts a function on the
stack, and each one returns at once. \texttt{setTimeout} doesn't wait zero
milliseconds; it hands its callback to libuv and returns. \texttt{then} doesn't
run its callback either: the promise is already resolved, so the callback goes
straight into the microtask queue. \texttt{process.nextTick} puts its callback
in Node's own queue. Only \texttt{console.log} prints anything. When the script
ends, the stack is empty, and Node empties the two small queues in a fixed
order: \texttt{nextTick} first (row~5), then the microtasks (row~6). That's why
\texttt{tick} comes before \texttt{promise}, although it was scheduled after
it. The timer callback is a task, and it has to wait for the loop to start
turning and reach the timers phase (row~7). The file is \texttt{order.js} on
purpose: Section~\ref{sec:modules} shows why the same lines in a \texttt{.mjs}
file print in another order.

The phases show up the same way: inside the callback of a file read, which runs
in poll, a \texttt{setImmediate} and a zero-millisecond \texttt{setTimeout} came
out immediate first in five runs out of five, because check comes right after
poll in the same turn and timers only in the next one. The cost in the concept
box is just as easy to provoke: a function that keeps chaining
\texttt{.then()} onto itself five million times kept the microtask queue from
ever emptying, and a timer set for zero milliseconds fired after two seconds.
Of the two small queues, \texttt{process.nextTick} is the older one, and Node's
own documentation now recommends \texttt{queueMicrotask()} instead.

With every queue in place, a real request can be followed through all of Node.
This program asks a server for something by name, \texttt{localhost}, not by
number, and prints the status code it gets back, \texttt{200}:

\begin{code}
async function get(url) {
  const res = await fetch(url);
  console.log(res.status);
}
get('http://localhost:8080/status');
\end{code}

Figure~\ref{fig:request-frames} follows it the same way as the four lines, with
two new columns: one for \texttt{get} while it's paused, and one for the work
that happens outside JavaScript.

\begin{figure}[htbp]
\centering
\begin{tikzpicture}[x=1cm, y=1cm]
  \tikzset{
    fr/.style={hf/base, draw=hfReader, fill=hfReader!10, font=\ttfamily\scriptsize,
      minimum width=24mm, minimum height=4.4mm, inner sep=1pt},
    me/.style={fr, line width=1.1pt, fill=hfReader!22},
    nd/.style={fr, draw=hfBus, fill=hfBus!10, font=\sffamily\scriptsize},
    pz/.style={fr, dashed, draw=black!45, fill=black!4, minimum width=20mm},
    mq/.style={fr, dashed, fill=hfReader!4, minimum width=21mm},
    when/.style={hf/note, anchor=west, align=left, text width=30mm},
    outside/.style={hf/note, anchor=west, align=left, text width=33mm, text=hfSource},
    head/.style={hf/note, font=\sffamily\footnotesize\bfseries, align=center},
    none/.style={hf/note, text=black!35},
    path/.style={-{Stealth[length=2mm]}, line width=0.9pt, draw=hfReader},
  }
  \def\rh{2.1}
  \node[head] at (4.6,1.45) {call stack};
  \node[head] at (7.05,1.45) {paused};
  \node[head, text=hfReader] at (9.35,1.45) {microtask\\queue};
  \node[head, text=hfSource] at (12.6,1.45) {outside JavaScript};
  \draw[black!12] (-0.1,1.05) -- (14.4,1.05);
  \foreach \r in {1,...,4} \draw[black!12] (-0.1,{-\rh*\r+1.05}) -- (14.4,{-\rh*\r+1.05});
  % 1
  \node[when] at (0,0) {\textbf{1.} \texttt{get} calls \texttt{fetch}, which starts the request and returns a promise that isn't resolved yet.};
  \node[fr] at (4.6,-0.47) {get.mjs};
  \node[me] (g1) at (4.6,0) {get};
  \node[fr] at (4.6,0.47) {fetch};
  \node[none] at (7.05,0) {---};
  \node[none] at (9.35,0) {empty};
  \node[outside] at (10.95,0) {nothing yet};
  % 2
  \node[when] at (0,{-1*\rh}) {\textbf{2.} \texttt{await} pauses \texttt{get}. The script ends, and no JavaScript runs at all.};
  \node[none] at (4.6,{-1*\rh}) {empty};
  \node[pz, draw=hfReader, text=hfReader] (g2) at (7.05,{-1*\rh}) {get};
  \node[none] at (9.35,{-1*\rh}) {empty};
  \node[outside] at (10.95,{-1*\rh}) {a pool thread asks the operating system for the address of \texttt{localhost}};
  % 3
  \node[when] at (0,{-2*\rh}) {\textbf{3.} The address arrives. The loop runs a piece of Node's code, a task, which opens the connection.};
  \node[nd] at (4.6,{-2*\rh}) {Node: connect};
  \node[pz, draw=hfReader, text=hfReader] (g3) at (7.05,{-2*\rh}) {get};
  \node[none] at (9.35,{-2*\rh}) {empty};
  \node[outside] at (10.95,{-2*\rh}) {the operating system sends the request and watches for the answer; no thread waits};
  % 4
  \node[when] at (0,{-3*\rh}) {\textbf{4.} The answer arrives. Another task reads it and resolves the promise.};
  \node[nd] at (4.6,{-3*\rh}) {Node: read the answer};
  \node[none] at (7.05,{-3*\rh}) {---};
  \node[mq] (g4) at (9.35,{-3*\rh}) {rest of get};
  \node[outside] at (10.95,{-3*\rh}) {done};
  % 5
  \node[when] at (0,{-4*\rh}) {\textbf{5.} The task ends. Node empties the microtask queue: \texttt{get} goes on after the \texttt{await} and prints \texttt{200}.};
  \node[me] (g5) at (4.6,{-4*\rh-0.23}) {get};
  \node[fr] at (4.6,{-4*\rh+0.24}) {console.log};
  \node[none] at (7.05,{-4*\rh}) {---};
  \node[none] at (9.35,{-4*\rh}) {empty};
  \node[outside] at (10.95,{-4*\rh}) {done};
  % the path of get
  \draw[path] (g1.east) to[out=0, in=90] (g2.north);
  \draw[path] (g2.south) -- (g3.north);
  \draw[path] (g3.south) to[out=-90, in=90] (g4.north);
  \draw[path] (g4.south) to[out=-90, in=0] (g5.east);
\end{tikzpicture}
\caption{The request from the code above, one row per moment. Follow
\texttt{get}, in blue: on the stack in row~1, paused in rows 2 and 3 while the
work happens outside JavaScript, sent to the microtask queue when the promise
resolves (row~4), and back on the stack in row~5. The purple boxes are Node's
own code, which the loop runs as tasks.}
\label{fig:request-frames}
\end{figure}

Row~2 is the one that surprises people: the request is in flight and no
JavaScript is running at all. \texttt{await} doesn't wait the way a blocking
call does. It takes \texttt{get} off the stack, keeps its variables aside with
the promise, and lets the script end. The waiting happens outside JavaScript,
in two places. The name goes to a pool thread, because the system call that
looks it up, \texttt{getaddrinfo}, can only block. The request goes to the
operating system, which watches the connection for free. Each time one of them
finishes, the loop wakes up and runs a piece of Node's own code as a task (rows
3 and 4). The second task resolves the promise, and that is the only thing that
brings \texttt{get} back.


\section{Two module systems}
\label{sec:modules}

\begin{history}{2009--2019 \textperiodcentered{} CommonJS and ES modules}
JavaScript was made in 1995 to add a few lines to a web page, and it had no way
for one file to use another. In the browser, scripts were loaded in order with
\texttt{<script>} tags and shared everything through global variables.

That was impossible on a server. In January 2009, Kevin Dangoor, then at
Mozilla, wrote a blog post called ``What Server Side JavaScript needs''. The
first item was a module system. The group that formed around it, first called
ServerJS and then \textbf{CommonJS}, specified \texttt{require()}: a function
that loads a file, runs it, and returns what it exported. Node adopted it the
same year.

Meanwhile the language itself was getting its own. Since 1997, JavaScript has
been defined by a standard, \textbf{ECMAScript}, kept by a committee of the
Ecma International standards body, TC39. It needed a different name because
``JavaScript'' was a trademark of Sun. Versions were numbered (ES3 in 1999, ES5
in 2009) until the sixth, in June 2015. That one was renamed \textbf{ES2015},
and since then a new version comes out every year, named after the year. It was
the largest change in the language's history, and among many other things it
added \texttt{import} and \texttt{export}, now called \textbf{ES modules}, or
ESM. Browsers adopted them. Node needed four more
years to make them work alongside a decade of \texttt{require}, and turned them
on without a flag in version 13.2, in November 2019.
\end{history}

The two systems look like two spellings of the same thing, but they work
differently. \texttt{require()} is an ordinary function call that loads a file synchronously,
in the middle of running code. An \texttt{import} is a declaration the runtime
reads before running anything, so it can find, fetch and link the whole graph
of files first, even asynchronously.

That difference is why \texttt{await} at the top of a file only works in ESM:

\begin{code}
// wait.cjs
const { setTimeout } = require('node:timers/promises');
await setTimeout(100);
console.log('ready');
\end{code}

\begin{terminal}
$ node wait.cjs
SyntaxError: await is only valid in async functions and the top level bodies of modules
$ node wait.mjs
ready
\end{terminal}

The \texttt{.mjs} file has the same three lines, written with \texttt{import}.
A CommonJS file can't wait, because whoever called \texttt{require()} expects
the result back immediately. An ES module can, because the runtime is already
loading the graph asynchronously (Figure~\ref{fig:cjs-esm}).

\begin{figure}[htbp]
\centering
\begin{tikzpicture}[x=0.92mm, y=1mm, bar/.style={font=\sffamily\scriptsize, text=white}]
  % --- CommonJS ---
  \node[hf/note, anchor=west, font=\sffamily\footnotesize\bfseries] at (-18,64) {CommonJS: \texttt{require()} runs the other file in the middle of a line};
  \node[hf/note, anchor=east] at (-2,54.5) {\texttt{main.js}};
  \node[hf/note, anchor=east] at (-2,44.5) {\texttt{config.js}};
  \fill[hfReader] (0,52) rectangle (30,57);
  \node[bar] at (15,54.5) {runs};
  \draw[black!35, dashed, line width=0.6pt] (30,54.5) -- (80,54.5);
  \node[hf/note] at (55,58.5) {stuck on \texttt{cfg = require(...)}};
  \draw[hf/ask] (30,52) -- (30,47);
  \fill[hfSource!75] (30,42) rectangle (80,47);
  \node[bar] at (55,44.5) {runs, top to bottom};
  \draw[hf/flow, draw=hfSource] (80,47) -- node[hf/note, right] {returns \texttt{cfg}} (80,52);
  \fill[hfReader] (80,52) rectangle (140,57);
  \node[bar] at (110,54.5) {uses \texttt{cfg}};
  \node[hf/note, text=hfAlert] at (55,37) {an \texttt{await} here would hand \texttt{main.js} a \texttt{cfg} that isn't ready};
  % --- ES modules ---
  \node[hf/note, anchor=west, font=\sffamily\footnotesize\bfseries] at (-18,26) {ES modules: every \texttt{import} is read first, then the files run in order};
  \node[hf/note, anchor=east] at (-2,16.5) {runtime};
  \node[hf/note, anchor=east] at (-2,8.5) {\texttt{config.js}};
  \node[hf/note, anchor=east] at (-2,0.5) {\texttt{main.js}};
  \fill[hfBus] (0,14) rectangle (30,19);
  \node[bar] at (15,16.5) {reads imports};
  \fill[hfSource!75] (30,6) rectangle (48,11);
  \node[bar] at (39,8.5) {runs};
  \draw[hfSource, dashed, line width=0.6pt] (48,6) rectangle (95,11);
  \node[hf/note] at (71.5,8.5) {\texttt{await}, 100 ms};
  \fill[hfSource!75] (95,6) rectangle (105,11);
  \fill[hfReader] (105,-2) rectangle (140,3);
  \node[bar] at (122.5,0.5) {runs, \texttt{cfg} ready};
  \node[hf/note] at (52,-6.5) {\texttt{main.js} hasn't started, so nobody is waiting on a half-done file};
\end{tikzpicture}
\caption{Why only an ES module can wait at the top. Above, \texttt{require()}
is a function call in the middle of \texttt{main.js}, which can't continue
until it returns. Below, the runtime reads every \texttt{import} before running
any code, so it can run \texttt{config.js} to the end, waits included, before
\texttt{main.js} starts.}
\label{fig:cjs-esm}
\end{figure}

Node decides which system a file uses by its extension (\texttt{.cjs} or
\texttt{.mjs}) or, for \texttt{.js}, by the \texttt{"type"} field of the nearest
\texttt{package.json}. Without that field, Node assumes CommonJS, for
compatibility with 2009. ESM also asks for the file extension in every import,
because it was designed for the browser, where an import is a URL and nobody
guesses extensions on a URL.

Loading asynchronously has one visible side effect. Saved as \texttt{order.mjs},
the four lines of Figure~\ref{fig:order-frames} print \texttt{script},
\texttt{promise}, \texttt{tick}, \texttt{timer}: the promise now comes before
the \texttt{nextTick}. An ES module's code runs as one step of that
asynchronous loading, inside a microtask. When the code ends, V8 is still
emptying the microtask queue it was running in, and Node only gets to its
\texttt{nextTick} queue afterwards. A CommonJS file runs as a plain task, so
\texttt{nextTick} goes first. Code that depends on the order between the two
small queues is a mistake in either kind of file.

CommonJS is legacy. A new project should be ESM from the first file and set
\texttt{"type": "module"}.

\section{Types that disappear}

\begin{history}{2012 \textperiodcentered{} Anders Hejlsberg}
By 2010, teams at Microsoft and elsewhere were writing applications in
JavaScript with hundreds of thousands of lines. The language gave them nothing
to hold on to: no declared types, so no editor that could say what a function
expected, and no way to rename something safely.

There were two kinds of answer. Google's GWT (2006) and Dart (2011) asked
programmers to write in another language and compile it to JavaScript. In
October 2012, Microsoft announced \textbf{TypeScript}, led by Anders Hejlsberg,
who had already designed Turbo Pascal, Delphi and C\#. TypeScript took the
opposite route: every JavaScript program is already a TypeScript program, and
the types you add are checked and then removed. The output is the JavaScript
you would have written by hand.
\end{history}

\begin{concept}{Type erasure}
TypeScript's types exist only while compiling. The compiler checks them and
then deletes them, and the program that runs has no trace of them.

\textbf{What it fixes:} adopting TypeScript costs nothing at runtime, and the
output runs anywhere JavaScript runs.

\textbf{What it costs:} nothing checks data that comes from outside the program
while it runs. A request body, a file or an environment variable has whatever
shape it has, regardless of what the types say.
\end{concept}

Compile \texttt{const port:\ number = 8080} and what comes out is
\texttt{const port = 8080}. The types are gone. What they were worth happened
before that, in the check.

One common slip shows how much that check is worth. Express gives its response
object a \texttt{send} method, and people who learned Express write
\texttt{res.send('hello')} out of habit. Node's own response object has no
\texttt{send}. It has \texttt{end}. Run without the compiler, the server starts,
says it's listening, and dies on the first request with ``res.send is not a
function'', which may happen hours later on another machine. With the compiler,
the program never starts: the check stops on that line and asks ``Did you mean
\texttt{end}?''.

The compiler knows Node's response object because someone described it in
TypeScript. Those descriptions come in a separate package, \texttt{@types/node},
from \textbf{DefinitelyTyped}, a community repository Boris Yankov started in
2012 for libraries that didn't ship their own types. They have to match the
Node that runs the program. With the types of a newer Node, the compiler
accepts functions the runtime doesn't have, and the check lies.

\section{The edit--run loop}

Compiling first and running second is safe and slow. Every change becomes
stop, compile, start, try. People don't work like that for long, and the tools
that fix it have their own lineage.

\begin{history}{2010--2025 \textperiodcentered{} Restarting on save}
In 2010, Remy Sharp wrote \textbf{nodemon}, which restarts a Node process
whenever a file changes. In 2015, Blake Embrey wrote \textbf{ts-node}, which
compiles TypeScript in memory so it can run without a build step. It
type-checked on every start, which made it slow. In 2020, Evan Wallace, a
co-founder of Figma, released \textbf{esbuild}, a JavaScript and TypeScript
compiler written in Go that is fast because it removes types without checking
them. In 2022, Hiroki Osame released \textbf{tsx}, which puts esbuild and
restart-on-save together.

Node then absorbed the idea. Version 22.6, in mid-2024, could strip types
behind a flag. Version 23.6, in January 2025, did it by default, and so did
22.18 that July. It only handles syntax that can be erased: TypeScript features
that generate code, like \texttt{enum}, don't pass.
\end{history}

\begin{figure}[htbp]
\centering
\begin{tikzpicture}
  \node[hf/source, minimum width=22mm] (src) {\texttt{.ts} file};
  \node[hf/bus, right=18mm of src, yshift=10mm, minimum width=32mm] (tsc) {\texttt{tsc}\\{\footnotesize checks, then writes}};
  \node[hf/source, right=10mm of tsc, minimum width=18mm] (js) {\texttt{.js}};
  \node[hf/reader, right=10mm of js, minimum width=18mm] (n1) {node};
  \node[hf/bus, right=18mm of src, yshift=-10mm, minimum width=32mm] (strip) {tsx, node\\{\footnotesize erases and runs}};
  \node[hf/reader, right=38mm of strip, minimum width=18mm] (n2) {running};
  \draw[hf/flow, black!60] (src.east) -- (tsc.west);
  \draw[hf/flow, black!60] (tsc) -- (js);
  \draw[hf/flow, black!60] (js) -- (n1);
  \draw[hf/flow, black!60] (src.east) -- (strip.west);
  \draw[hf/flow, black!60] (strip) -- node[hf/note, text=hfAlert, above] {no type check} (n2);
\end{tikzpicture}
\caption{Two ways from TypeScript to a running program. The fast one is fast
because it skips the check.}
\end{figure}

The fast loop gives the misspelled method its old behavior back: the process
restarts, the server comes up, and the error waits for a request. The build
would still catch it, but nobody runs a build every time they save. So the
check has to run somewhere else, as its own command, \texttt{tsc} with
\texttt{--noEmit}, every time the code is about to be shared. A check that
only runs when somebody remembers to run it doesn't count as a check.

\section{What is actually installed}

\begin{history}{2010--2017 \textperiodcentered{} npm and the lock file}
Isaac Schlueter released npm in 2010, and it became Node's package manager.
Packages declared their dependencies as version ranges, following Semantic
Versioning, which Tom Preston-Werner, a co-founder of GitHub, had written down
around the same time: the first number changes when something breaks, the
second when something is added, the third for fixes. A range like
\texttt{\textasciicircum 2.1.0} means ``2.1.0 or anything newer before 3''.

Ranges meant that two installs of the same project on different days could get
different code, because a dependency of a dependency had published a new
version in between. In October 2016, Facebook, with Google and others, released
\textbf{Yarn}, whose main promise was a lock file: the exact version of every
package in the tree, written down, so every machine installs the same thing.
npm adopted the idea in version 5, in 2017, with \texttt{package-lock.json}.
\end{history}

\begin{figure}[htbp]
\centering
\begin{tikzpicture}
  \node[hf/source, minimum width=36mm, minimum height=16mm] (pkg) {\texttt{package.json}\\{\footnotesize \texttt{some-lib \textasciicircum 2.1.0}}\\{\footnotesize what you accept}};
  \node[hf/bus, right=22mm of pkg, minimum width=36mm, minimum height=16mm] (lock) {\texttt{package-lock.json}\\{\footnotesize \texttt{some-lib 2.4.3} + hash}\\{\footnotesize what was chosen}};
  \node[hf/store, right=22mm of lock, minimum width=30mm, minimum height=16mm] (nm) {\texttt{node\_modules}\\{\footnotesize what is on disk}};
  \draw[hf/flow, black!60] (pkg) -- node[hf/note, above] {npm install} (lock);
  \draw[hf/save] (lock) -- node[hf/note, above] {npm ci} (nm);
  \draw[hf/pain, dashed, -{Stealth[length=2mm]}] ([yshift=-10mm]pkg.south) -- ([yshift=-10mm]nm.south)
    node[pos=0.5, hf/note, text=hfAlert, below=1.5mm] {editing \texttt{package.json} by hand changes neither of the other two};
\end{tikzpicture}
\caption{Three different things. Only \texttt{npm install} brings them into
agreement.}
\end{figure}

\begin{concept}{Lock file}
A file, generated by the package manager and committed with the code, that
records the exact version and content hash of every package the project
installed, including the dependencies of dependencies.

\textbf{What it fixes:} the code that ran on your machine is the code that runs
on the server, and a new release of a package you've never heard of can't
change your program overnight.

\textbf{What it costs:} updates stop happening by themselves. Somebody has to
decide to update, and review a lock diff that can be hundreds of lines for a
single new dependency.
\end{concept}

Those hundreds of lines are often honest. A tool written in Go, like esbuild,
is published as one package per operating system and processor, all listed as
optional dependencies; the lock records every one of them, and npm downloads
only the one for your machine.

The \texttt{package.json} also separates what the program needs to run
(\texttt{dependencies}) from what it only needs to be built and checked
(\texttt{devDependencies}). The compiler, the type descriptions and the test
runner are the second kind. The image that runs in production installs only the
first kind, which makes it smaller and leaves less code there to attack.

Commit the lock file, and install from it with \texttt{npm ci} everywhere except
your own machine. A project without a lock file installs a slightly different
program every week.

\section{A framework}

\begin{history}{2007--2024 \textperiodcentered{} Sinatra, Express, Fastify}
In 2007, Blake Mizerany released Sinatra for Ruby: a web application as a list
of routes, each one a method, a path and a block of code. In 2010, TJ
Holowaychuk brought that style to Node with \textbf{Express}, built on a chain
of \emph{middleware}: functions that each receive the request, do something,
and pass it on. Express became what Node servers were written in. Its version 4
came out in 2014, and version 5 only in 2024. Until 5, an error thrown inside an
\texttt{async} handler wasn't caught, and the request simply hung.

In 2016, Matteo Collina and Tomas Della Vedova started \textbf{Fastify}, which
reached 1.0 in March 2018. It came from measuring what Express spent per
request: the middleware chain, and a generic \texttt{JSON.stringify} for every
response. Fastify asks each route to declare a schema for what it receives and
returns, and compiles that schema into a function that validates or serializes
exactly that shape.
\end{history}

Go back to the first server in this chapter, the one written on the bare
\texttt{http} module. It had one function that received every request and
checked, with an \texttt{if}, whether the method was \texttt{GET} and the
address was exactly \texttt{/status}. Everything else got a 404. With a
framework, the program turns around: you write one function per route, and the
framework decides which one to call. Here is the same \texttt{/status} route in
Fastify, plus a second route that receives a number:

\begin{code}
app.get('/status', async () => ({ ok: true }));

app.post('/speed', {
  schema: {
    body: {
      type: 'object',
      required: ['factor'],
      properties: { factor: { type: 'number', minimum: 0.1, maximum: 10 } },
    },
  },
}, async (req) => ({ accepted: req.body }));
\end{code}

Compare it with the bare version. The \texttt{if} is gone: neither function looks
at the method or the address, because \texttt{app.get} and \texttt{app.post}
already said which requests each one is for. And the \texttt{/speed} route
declares what its body must be, an object with a field \texttt{factor} that is
a number between 0.1 and 10, while the function itself never checks it.

To try the routes, I didn't start the server. Fastify has a method,
\texttt{inject}, that builds a request in memory and runs it through the
application as if it had arrived over the network: the same routing, the same
schema check, the same function. Nothing listens on a port and nothing crosses
the network. It isn't a test by itself, but it's what tests are built on: a test
can check a route without starting a server, finding a free port and waiting
for it to come up. These are four requests sent that way, and what the bare
server did with the same ones:

\begin{itemize}
  \item \texttt{GET /status?x=1} got \textbf{200} and \texttt{\{"ok":true\}}.
        The bare server answered 404, because it compared the whole address,
        query string included, as text. Fastify separates the path from the
        query before choosing the route.
  \item \texttt{POST /speed} with \texttt{factor} 2 got \textbf{200}, and the
        function received the body already parsed. The bare server would have
        needed code to collect the body from the stream and parse it.
  \item \texttt{POST /speed} with \texttt{factor} 50 got \textbf{400} and the
        message ``body/factor must be <= 10''. The function never ran: the
        schema turned the request away and said why. In the bare version, that
        check would be more \texttt{if}s inside the function.
  \item \texttt{POST /status} got \textbf{404}, ``Route POST:/status not
        found''. The bare server also answered 404.
\end{itemize}

The last one is wrong in both. The path exists, only not for that method, and HTTP has
had a status for that since 1997: \emph{405 Method Not Allowed}, which comes
with a header listing the methods the path does accept. Neither Express nor
Fastify sends it by default. It's a small mistake, and both make it.

For a new TypeScript server, Express is the wrong default today. Its types live
in DefinitelyTyped, maintained apart from the code, and it took until 2024 to
handle errors in \texttt{async} handlers. Fastify ships its own types, validates
input from the route's schema, and can be tested without a network.

\section{Settings from outside}

The same program has to run on a laptop, in a test, and in a cluster, with a
different port, address and log level in each. The settings can't be in the
code.

\begin{history}{2004--2019 \textperiodcentered{} Failing early}
In 2004, Jim Shore wrote ``Fail Fast'' for \emph{IEEE Software}: a program that
detects a problem should stop immediately and say so, instead of carrying on
with a bad value and failing later, far from the cause. In 2011, Adam Wiggins,
a co-founder of Heroku, published \textbf{The Twelve-Factor App}, a list of
rules for programs that run on platforms like his. Rule three: configuration
lives in environment variables, so the same build runs everywhere and no
password ends up in the repository. In 2019, Alexis King wrote ``Parse, don't
validate'', which gave a name to the right way of checking input: turn it, once,
at the edge, into a type that can't hold an invalid value.
\end{history}

Environment variables come with a problem the types don't see. Every one of
them is a string or missing, and JavaScript is happy to add a number to a string:

\begin{terminal}
$ TIMEOUT_MS=5000 node -e "console.log(process.env.TIMEOUT_MS + 1)"
50001
\end{terminal}

The fix, after type erasure, has to happen while the program runs. Libraries
for it go back a while. Joi, from the hapi framework, around 2012, checked
values at runtime but knew nothing about TypeScript. io-ts, by Giulio Canti in
2017, derived the type from the check, in a heavily functional style. Colin
McDonnell's \textbf{zod}, in 2020, did the same with an API that reads like
ordinary code, and became the common choice.

\begin{code}
import { z } from 'zod';

const Env = z.object({
  TIMEOUT_MS: z.coerce.number().int().positive().default(5000),
  MODE: z.enum(['live', 'replay']).default('live'),
});

const result = Env.safeParse(process.env);
if (!result.success) {
  console.error(z.prettifyError(result.error));
  process.exit(1);
}
const env = result.data; // { TIMEOUT_MS: number; MODE: 'live' | 'replay' }
console.log(env.TIMEOUT_MS + 1, env.MODE);
\end{code}

Without the variable, the program prints \texttt{5001 live}. With
\texttt{TIMEOUT\_MS=5s}, it doesn't start:

\begin{terminal}
$ TIMEOUT_MS=5s node env.ts
(*$\times$*) Invalid input: expected number, received NaN
  (*$\rightarrow$*) at TIMEOUT_MS
\end{terminal}

\begin{concept}{Parse at the edge}
Check everything that comes from outside the program once, where it enters,
and turn it into a typed value. From then on, the code uses the typed value and
never looks at the raw input again.

\textbf{What it fixes:} a bad setting stops the program at startup, with a
message naming it, instead of breaking the first operation that uses it, an
hour later.

\textbf{What it costs:} every input needs a schema, written and kept in step
with the code that uses it.
\end{concept}

An empty \texttt{TIMEOUT\_MS=} is refused too, and that one surprises people.
It doesn't fall back to the default, because the variable is defined, just
empty, and an empty string becomes the number zero, which isn't positive.
Failing there is correct. An empty variable in a deployment is almost always a
mistake in the deployment.

\section{What the compiler doesn't check}

\begin{history}{1978 \textperiodcentered{} Stephen C.\ Johnson}
At Bell Labs, Stephen C.\ Johnson wrote \textbf{lint}, a program that read C
code and complained about things the compiler accepted but that were probably
wrong: variables used before being set, functions called with the wrong
arguments, code that would only work on one machine. He kept it separate from
the compiler on purpose. The compiler's job was to translate quickly; lint's
was to be suspicious. In his 1978 paper he wrote that ``the general notion of
having two programs is a good one''.

JavaScript got its first lint in 2002, Douglas Crockford's JSLint, which
enforced Crockford's own opinions with no way to turn them off. Anton
Kovalyov forked it in 2011 as JSHint, which was configurable. In 2013, Nicholas
Zakas released \textbf{ESLint}, in which every rule is a plugin. That design
is why TypeScript got its own lint: in 2019, the separate TSLint project was
retired, and typescript-eslint, a set of ESLint plugins that can ask the
TypeScript compiler about types, took its place.
\end{history}

Some mistakes are well-typed. Calling an \texttt{async} function without
waiting for it is valid code, with a valid type:

\begin{code}
async function save(value: number): Promise<void> {
  await new Promise((resolve) => setTimeout(resolve, 100));
  if (value < 0) throw new Error('negative');
  console.log('saved', value);
}

save(42);
console.log('done');
\end{code}

It prints \texttt{done} and then \texttt{saved 42}: the program announces
it's finished before the save happens, and if the save fails, nobody is waiting
to hear it. The compiler accepts this file without a word, because every type
in it is correct. The mistake is in how the value is used. A lint rule that can
see types catches it: typescript-eslint's \texttt{no-floating-promises} flags
the line \texttt{save(42)} and asks for an \texttt{await}, a \texttt{.catch}, or
an explicit \texttt{void} to say the result is being ignored on purpose.

To see types, the linter calls the compiler's own JavaScript API. That
dependency surfaced in 2025. In March, Hejlsberg announced that TypeScript was
being ported to Go, for compile times around ten times shorter. TypeScript 6.0,
in March 2026, was the last version written in JavaScript, and 7.0, the Go
port, came out that July without a compatible version of the API that
typescript-eslint uses.

\begin{question}{The TypeScript compiler is written in Go? Do I need Go now?}
No. A compiler is a program like any other. It reads text and writes text, and
nothing requires it to be written in the language it reads. The first ones
couldn't be: IBM's FORTRAN compiler, in 1957, was written in assembly, because
there was no FORTRAN yet to write it in.

Go is needed to build the compiler, not to run it, and what Microsoft publishes
is the result: an executable in the processor's own instructions. npm delivers
it the same way it delivers esbuild, as described earlier in this chapter. On a
Windows machine, installing \texttt{typescript} 7.0.2 brings a package that
lists more than a dozen optional packages, one per operating system and
processor, and npm downloads only \texttt{@typescript/typescript-win32-x64}.
Inside it is \texttt{tsc.exe}, a 24~MB Windows executable. The \texttt{tsc}
command a project runs is still a JavaScript file, two lines long, whose only
job is to start it.
\end{question}

Writing a compiler in the language it compiles came later, and it was a
milestone.

\begin{history}{1962 \textperiodcentered{} A compiler that compiles itself}
At MIT in 1962, Tim Hart and Mike Levin wrote a compiler for Lisp in Lisp. To
get the first copy running, they ran its source in the Lisp interpreter that
already existed and had it compile itself. From then on, the compiled compiler
compiled its own new versions. A compiler written in the language it compiles
is called \textbf{self-hosting}, and using an older tool to produce the first
copy is called \textbf{bootstrapping}. Dennis Ritchie's C compiler at Bell Labs
was written in C from its early versions in the 1970s. Being self-hosting
became a sign that a language was serious: if it was good enough for its own
compiler, it was good enough for most things.
\end{history}

TypeScript was self-hosting from 2012 until 2026: written in TypeScript,
compiled to JavaScript by an earlier version of itself, and run by Node. That
had a cost. Node runs code on one thread, under an event loop designed for
programs that mostly wait on the network, and a type checker doesn't wait. It
computes, for a long time, and on a large project it was slow. Hejlsberg's team
chose Go over Rust because they wanted to port the existing code line by line,
not redesign it, and Go, with its garbage collector, let them keep the same
structure. Go compiles to machine code and uses every processor core at once.

That is also why the linter broke. While the compiler was JavaScript, it ran in
the same process as the linter, which could call its functions and read its
objects directly, the way any library is used. A Go executable is another
program, in another process. The only way to use it is to send it messages and
read its answers, and that interface wasn't ready when 7.0 came out. A project
that lints with types stays on 6.0 until it is. A compiler is also a library
that other tools are built on, and making it faster broke them.

\section{Format, line endings, tests}

\begin{history}{2009--2017 \textperiodcentered{} The end of style arguments}
When Go was released in 2009, it came with \texttt{gofmt}, which rewrites any
Go file into the one official layout, and Go programmers stopped arguing about
where braces go. In January 2017, James Long released \textbf{Prettier}, which
did the same for JavaScript: it parses the code, throws away the original
layout, and prints the syntax tree again following an algorithm Philip Wadler
had published as ``A prettier printer''.
\end{history}

Arguing about formatting is useless. A team should pick a formatter, accept
its output, and spend the time on something else. Formatters and linters used
to fight over the same rules, so the usual setup turns off every formatting
rule in the linter and leaves layout to the formatter.

\begin{history}{1963--1981 \textperiodcentered{} Two motions of a teleprinter}
On a Teletype, ending a line took two mechanical motions: the
\emph{carriage return} (CR) sent the print head back to the left, and the
\emph{line feed} (LF) rolled the paper up one line. ASCII, in 1963, gave each
one a code, and a line ended with both. Multics and then Unix decided that a
line inside a file ends with LF alone, and the terminal driver adds the CR when
printing. CP/M, in the 1970s, kept both, and MS-DOS, in 1981, inherited them
from CP/M. Windows still ends lines with CR LF.
\end{history}

The consequence shows up in a repository edited on Windows and run on Linux.
Git for Windows, by default, converts line endings to CR LF when it checks
files out. A formatter that writes LF then reports every file as wrong. The fix
is a \texttt{.gitattributes} file that tells Git the repository's files use LF
on every machine, which is also what the Linux container will read.

Finally, tests. Facebook released Jest in 2014, and for years it was what
JavaScript projects tested with, but it needed extra compilers for TypeScript
and its support for ES modules stayed experimental. Vitest, from Anthony Fu
and the Vite team in December 2021, runs TypeScript and ES modules directly. A
test of an HTTP route doesn't need a network: if the code that \emph{builds}
the application is separate from the code that \emph{starts} it, the test
builds the application and sends it requests with \texttt{inject}, and no port
is ever opened.

\section{What it costs}

After all of this, the server answers one route, the same one the bare
\texttt{http} module answered with a dozen lines. In exchange, a misspelled
method fails the type check, a bad setting stops the startup, a forgotten
\texttt{await} fails the lint and a broken route fails a test, all before
anyone opens a browser.

The price is a dozen tools, each with its own configuration file and its own
release schedule, and the versions have to agree: the type descriptions with
the runtime, the linter with the compiler, the formatter with the line endings.
When one of them moves, as the compiler did in 2026, something downstream
breaks. Most of the work of keeping a Node project healthy is keeping those
versions in agreement.

\section*{Further reading}

\begin{itemize}
  \item D.\ Kegel, ``The C10K problem'', 1999,
        \url{http://www.kegel.com/c10k.html}.
  \item R.\ Dahl, ``Node.js'', JSConf EU, Berlin, 2009.
  \item K.\ Dangoor, ``What Server Side JavaScript needs'', 2009,
        \url{https://blueskyonmars.com/2009/01/29/what-server-side-javascript-needs}.
  \item S.~C.\ Johnson, ``Lint, a C Program Checker'', Bell Laboratories, 1978.
  \item J.\ Shore, ``Fail Fast'', \emph{IEEE Software} 21(5), 2004.
  \item A.\ Wiggins, \emph{The Twelve-Factor App}, 2011,
        \url{https://12factor.net}.
  \item A.\ King, ``Parse, don't validate'', 2019,
        \url{https://lexi-lambda.github.io/blog/2019/11/05/parse-don-t-validate/}.
  \item P.\ Wadler, ``A prettier printer'', in \emph{The Fun of Programming},
        Palgrave, 2003.
  \item S.\ McKenzie, C.\ Pojer and J.\ Kyle, ``Yarn: A new package manager for
        JavaScript'', Facebook Engineering, 2016,
        \url{https://engineering.fb.com/2016/10/11/web/yarn-a-new-package-manager-for-javascript/}.
\end{itemize}
